\documentclass[]{article}

\usepackage{amsmath}
\usepackage{multirow}
\usepackage{natbib}
\usepackage{graphicx} 


\begin{document}

The thermal Sunyaev-Zel'dovich (tSZ) effect, arising from the inverse Compton scattering of cosmic microwave background photons off hot, ionized gas, offers a direct probe of the integrated electron pressure along the line of sight. This makes it an indispensable tool for mapping the distribution of baryons in the universe, studying the thermodynamics of the intracluster and circumgalactic medium, and constraining models of galaxy formation. In particular, the morphology of the tSZ signal on small angular scales, from one to five arcminutes, encodes crucial information about the impact of energetic feedback from supernovae and active galactic nuclei. Upcoming experiments, such as the Simons Observatory, are set to map the microwave sky with the sensitivity and resolution required to access this information, but realizing this potential requires novel methods to extract the faint tSZ signal from complex observational data.

A formidable challenge in this endeavor is the contamination from other astrophysical signals and instrumental noise. At the small angular scales most relevant for feedback studies, the dominant foreground is the Cosmic Infrared Background (CIB), the cumulative emission from dusty, star-forming galaxies across cosmic time. The difficulty in separating the tSZ signal from the CIB is twofold. First, the CIB is a significant source of emission in the microwave bands. Second, and more critically, it is spatially correlated with the tSZ signal, as the star formation it traces is physically linked to the same feedback processes that heat the surrounding gas. This complex, non-Gaussian relationship violates the statistical assumptions underpinning traditional linear component separation methods, such as constrained Internal Linear Combination or Wiener Filtering. These techniques, which rely primarily on the distinct frequency spectra of the components, struggle to disentangle spatially correlated signals, often resulting in biased reconstructions with significant signal loss or residual foreground contamination.

In this paper, we reframe the tSZ map reconstruction as a guided super-resolution denoising problem and introduce a deep learning framework to address the limitations of linear methods. Instead of treating the CIB as a contaminant to be removed, our approach leverages the high-frequency observational channels where it is brightest as a spatial guide to reconstruct the underlying tSZ pressure map. Our framework consists of a two-stage model trained on simulated multi-frequency observations from the FLAMINGO simulation, designed to mimic the data from the Simons Observatory. The first stage is a Super-Resolution Denoising Autoencoder that learns to map the noisy, multi-frequency input to a clean, high-resolution (1-arcmin) tSZ map. The network is trained with a composite loss function that enforces both pixel-level accuracy and fidelity to the physical tSZ power spectrum, which prevents the generation of unphysical structures and ensures the statistical properties of the reconstruction are correct.

Recognizing that a single point-estimate map is insufficient for rigorous scientific analysis, the second stage of our framework transitions this deterministic model into a probabilistic one. We employ a Conditional Diffusion Model, which learns the full conditional probability distribution of the tSZ map given the observational data. By sampling from this distribution, we can produce not only a high-fidelity reconstruction but also generate robust, pixel-level uncertainty maps that quantify the model's confidence. We demonstrate that our framework significantly outperforms standard linear methods in recovering the tSZ signal, especially on the 1--5 arcminute scales. The resulting maps yield a tighter tSZ signal-mass scaling relation for galaxy clusters, and we show that the derived uncertainties are well-calibrated, providing a reliable measure of map fidelity for future cosmological and astrophysical analyses with the next generation of microwave surveys.

\end{document}
                